Grouped scanning adds blocks within each query, then merges survivors \cite{johnson_gpu}. It preserves computed-primary selection, not exact raw-vector neighbours.

\subsection{Candidate partition and deterministic ordering}

For a fixed query, concatenate the selected inverted lists in probe order
to obtain $n_c$ candidate occurrences. Assume that each visited database
row occurs once and that the index stays fixed during search. For
$\gscan\in\{1,2,4,8\}$, block $(q,j)$ owns the interval
\begin{equation}
 U_j=\left[\left\lfloor\frac{j n_c}{\gscan}\right\rfloor,
              \left\lfloor\frac{(j+1)n_c}{\gscan}\right\rfloor\right),
 \qquad j=0,\ldots,\gscan-1.
 \label{eq:5-partition}
\end{equation}
Here an interval denotes the corresponding candidates, rather than their
database IDs. The device computes its endpoints from the routed candidate
count. Groups differ in size by at most one candidate; a long list can
cross a group boundary. This balances candidate counts, whereas assigning
equal numbers of lists need not do so. Each thread finds its first list
through the probe-prefix offsets and advances its list cursor while
scanning its interval. Residual codes are not read at this stage.

The selection key is
\begin{equation}
 K_q(x)=\bigl(\primarydist(q,x),p(x)\bigr),
 \qquad
 (a,p)\prec(b,r)\iff a<b\ \text{or}\ (a=b\ \text{and}\ p<r),
 \label{eq:5-primary-key}
\end{equation}
where $p(x)$ is the physical row position in the fixed inverted-list
storage. It is unique among the query's candidates and independent of
the order of atomic appends. It is not the database ID used to break
final-result ties. With non-NaN scores, this defines a total order on
distinct candidates. All local comparisons and merge decisions use that
same order. Partition-independent selection additionally requires that
each candidate receive the same computed score in every compared path;
the ordering rule alone cannot remove floating-point scoring differences.

\subsection{Exact local retention and global coverage}

$\operatorname{Top}_c(V)$ is the first $\min(c,|V|)$ elements under $K_q$, with $c=\max\{k,\lceil\alpha k\rceil\}$. Each group needs its full local top-$c$: all global winners could lie in one group. A $c/\gscan$ quota can lose them.

\begin{resultblock}
\begin{lemma}[Local-buffer retention]
\label{lem:5-buffer}
Let $V_t$ be the candidates scored by a block through chunk $t$, and
$B_t$ its valid buffer after any compaction or filtering at that point.
Assume fixed complete keys, sufficient append capacity and exact
top-$c$ compaction. Let $\beta_t$ be the minimum of the complete shared
bounds observed so far, including the initial top sentinel. Then
\begin{equation}
 \operatorname{Top}_c(V_t)\cap\{x:K_q(x)\preceq\beta_t\}
 \subseteq B_t.
 \label{eq:5-buffer-invariant}
\end{equation}
With sharing disabled, the bound is the sentinel and the buffer contains
the whole prefix top-$c$.
\end{lemma}
\begin{proof}
The empty prefix satisfies the invariant. Local rejection has $c$ earlier witnesses, so the rejected key cannot become a later prefix winner. Shared rejection excludes only keys above $\beta_t$. Every other eligible winner is already retained or appended. Exact compaction keeps it because fewer than $c$ prefix keys precede it; filtering removes only keys above an observed bound. Induction over these operations proves the claim. The running minimum handles stale reads of previously published bounds.
\end{proof}
\end{resultblock}

\begin{resultblock}
\begin{proposition}[Coverage and truncated merging]
\label{prop:5-local-coverage}
For fixed keys under a common total order, let $U$ be the visited
universe and $V$ any retained subset. If
$\operatorname{Top}_c(U)\subseteq V\subseteq U$, then
$\operatorname{Top}_c(V)=\operatorname{Top}_c(U)$.
Pairwise sorted merging with truncation to $c$ preserves this result.
In particular, for disjoint groups,
\begin{equation}
 \operatorname{Top}_c\!\left(\bigcup_j U_j\right)
 =\operatorname{Top}_c\!\left(\bigcup_j
                    \operatorname{Top}_c(U_j)\right).
 \label{eq:5-topc-union}
\end{equation}
The statement includes universes smaller than $c$.
\end{proposition}
\begin{proof}
Every global winner is in $V$, and every other member of $V$ follows
the global winners. Selection from $V$ therefore recovers them.
For the independent-group special case, let $x$ belong to the global
top-$c$ and to group $U_j$. Fewer than $c$
candidates precede $x$ globally, so fewer than $c$ precede it within $U_j$.
Consequently $x\in\operatorname{Top}_c(U_j)$, and the local union contains
every global winner. Selecting its first $c$ therefore recovers precisely
the global winners in the same order. If the whole universe has fewer
than $c$ members, every member is retained.

For merging, the first $c$ of two sorted lists are their union's
top-$c$. A discarded candidate already has at least $c$ predecessors
within that union and cannot enter the top-$c$ of any larger union.
Induction over the merge tree establishes the claim.
\end{proof}
\end{resultblock}

\begin{figure}[tbp]
\centering
% Schematic only: illustrative keys, not experimental observations.
\begin{tikzpicture}[
  font=\small\linespread{1}\selectfont,
  box/.style={draw=ThesisNavy!65,rounded corners=2pt,align=center,
    inner sep=5pt,text width=2.65cm,minimum height=2.05cm},
  merged/.style={box,text width=5.55cm,minimum height=1.05cm,
    fill=ThesisNavy!5},
  arr/.style={-{Latex[length=2mm]},draw=ThesisNavy!80,thick}]
\node[box] (g0) at (0,0) {$\operatorname{Top}_3(U_0)$\\[2pt]
  $\mathbf{(1,8)}$\\$\mathbf{(2,6)}$\\$(2,9)$};
\node[box] (g1) at (3.5,0) {$\operatorname{Top}_3(U_1)$\\[2pt]
  $(4,1)$\\$(5,4)$\\\textit{length 2}};
\node[box] (g2) at (7,0) {$\operatorname{Top}_3(U_2)$\\[2pt]
  $\varnothing$\\\textit{length 0}};
\node[box] (g3) at (10.5,0) {$\operatorname{Top}_3(U_3)$\\[2pt]
  $\mathbf{(2,3)}$\\$(7,2)$\\$(8,5)$};
\node[merged] (m0) at (1.75,-2.5) {Merge and retain first 3\\[2pt]
  $\mathbf{(1,8)},\ \mathbf{(2,6)},\ (2,9)$};
\node[merged] (m1) at (8.75,-2.5) {Merge and retain first 3\\[2pt]
  $\mathbf{(2,3)},\ (7,2),\ (8,5)$};
\node[merged,text width=6.8cm,fill=ThesisNavy!12] (out) at (5.25,-4.3)
  {Global primary top-3\\[2pt]
   $\mathbf{(1,8)},\ \mathbf{(2,3)},\ \mathbf{(2,6)}$};
\draw[arr] (g0.south) -- (m0.north west);
\draw[arr] (g1.south) -- (m0.north east);
\draw[arr] (g2.south) -- (m1.north west);
\draw[arr] (g3.south) -- (m1.north east);
\draw[arr] (m0.south) -- (out.north west);
\draw[arr] (m1.south) -- (out.north east);
\end{tikzpicture}

\caption[Local coverage and truncated pairwise merging]{Local coverage and truncated pairwise merging, illustrated with
$c=3$ and $\gscan=4$. Entries are $(\primarydist,p)$ keys, not measured data.
Bold entries are the final global winners. A short group contributes its
valid prefix; an empty group contributes no entries. Every merge uses the
same key order and retains at most $c$ entries.}
\label{fig:5-global-selection}
\end{figure}

The CUDA implementation maintains a shared-memory buffer of $C_{\rm buf}$ key
pairs, where $C_{\rm buf}$ is a power of two and $C_{\rm buf}\geq c+T_{\rm blk}$ for a scan block of
$T_{\rm blk}$ threads. A compaction sorts this buffer with a bitonic network,
retains at most $c$ valid entries and records the last retained key as
the local threshold when $c$ entries exist. Before that point the
threshold is a sentinel above every valid key. Between compactions,
only candidates preceding the threshold are appended. Its $c$ witness
candidates make the threshold a safe upper bound for the best $c$ of
the scanned prefix, even as further candidates arrive.

Each chunk appends at most $T_{\rm blk}$ entries. Compaction precedes
any chunk that would exceed capacity, leaving at most $c$ entries and
space for the next chunk because $C_{\rm buf}\geq c+T_{\rm blk}$.
This supplies the capacity premise of \Cref{lem:5-buffer} without a
per-thread quota. With sharing, \Cref{prop:5-shared-bound} bounds every
observed threshold above the global $c$th key. The buffer lemma and
coverage proposition therefore preserve global winners through filtering
and merging.

Materialising all $n_c$ score/position pairs needs $O(n_c)$ writes/storage. Fused scoring instead repeatedly sorts bounded local buffers and emits $O(\gscan c)$ entries. Register selection is another established option \cite{johnson_gpu}; the experiments test the chosen scan variants.

\subsection{Merge implementation, storage and boundary cases}

For $\gscan>1$, one 256-thread block merges each query over $\log_2\gscan$ rounds. Each pair emits $\min(c,a+b)$ entries. Merge Path finds thread slices by binary search, then each thread merges its slice \cite{gpu_merge_path}.
All comparisons use \cref{eq:5-primary-key}; block barriers separate
rounds. \Cref{fig:5-global-selection} illustrates the truncation invariant.
The tree emits at most $c(\gscan-1)$ key pairs across its internal nodes,
in addition to the partition searches and synchronisation work.

The two global score/position workspaces have capacities
$[b_{\rm cap},\gscan,c]$ and $[b_{\rm cap},\gscan/2,c]$ and alternate between
rounds. With four-byte scores and positions, their combined payload is
\begin{equation}
 S_{\rm pairs}=8b_{\rm cap}c\left(\gscan+\frac{\gscan}{2}\right)
 \quad\text{bytes},\qquad \gscan>1.
 \label{eq:5-group-workspace}
\end{equation}
Group counts require another $4b_{\rm cap}\gscan$ bytes; shared bounds add
$8b_{\rm cap}$ when provisioned. These are additional to final survivor
arrays and other search workspace. Allocation grows before graph capture
and is retained for reuse, so choosing a smaller $\gscan$ later does not
necessarily release the reserved memory. The merge keeps only list
lengths in shared memory rather than all $\gscan c$ keys.

Valid counts delimit each local prefix. Sort padding uses $(+\infty,\mathrm{INT\_MAX})$; valid positions must be smaller. Empty groups have count zero. Output padding instead uses $(+\infty,-1)$, which refinement skips and merging must not treat as valid. $\gscan=1$ writes survivors directly, without merging.

$\gscan=0$ labels one block per query. Its default selector is exact by distance but can choose different cutoff-tie IDs from the complete-key path. The lossy \texttt{K\_LOCAL=4} alternative requires \texttt{JHQ\_EXACT\_TOPCK=0}; paper-faithful builds disable it.

For distance-only compaction, every rejected score has $c$ retained witnesses no larger than it. The final buffer therefore contains all scores below the cutoff and enough equal scores, without fixing tied IDs. T4 checks this preorder and complete-key grouped selection (\cref{sec:7-strict-selection}). The reference sorts GPU-produced scores, rather than recomputing original-space distances on the CPU.